%%%%%%%%%%%%%%%%%%%%%%%%%%%%%%%%%%%%%%%%%%%%%%%%%%%%%%%%%%%%%%%%%%%%%%%%%%%%%%%%
\chapter{Syntax\label{chap:Syntax}}
%%%%%%%%%%%%%%%%%%%%%%%%%%%%%%%%%%%%%%%%%%%%%%%%%%%%%%%%%%%%%%%%%%%%%%%%%%%%%%%%

This chapter defines the grammar of \ASLvOneterm{}.
The grammar is presented as a context-free grammar~\cite{HU79} with the help of
two shorthands --- \emph{inlined derivations} and \emph{parametric productions} ---
inspired  by the Menhir Parser Generator~\cite{MenhirManual} for the
\href{https://ocaml.org/}{OCaml language}.
These shorthands can be considered as macros over standard context-free grammar
notation, which can be expanded away to yield a standard context-free grammar.

\section*{Ambiguities and LR(1) grammars}

In most cases, \ASLvOneterm{}'s context-free grammar is unambiguous.
This means that for any given text, either:
\begin{itemize}
  \item It cannot be derived from the grammar. The text is therefore invalid
        ASL syntax, and causes a parse error. An example is \texttt{1 + * 3},
        which cannot be derived
        from the grammar in \secref{ASLGrammar} (intuitively, because \texttt{+}
        and \texttt{*} are binary operators that require two sub-expressions).
        See \ExampleRef{Failing to Parse an Invalid Expression}.
  \item There is exactly one way to derive it from the grammar. The text is
        valid ASL syntax, and the derivation from the grammar describes how it
        should be bracketed.
\end{itemize}

In some cases, there are ambiguities in \ASLvOneterm{}'s context-free grammar.
For example, the expression \texttt{1 + 2 * 3} can be derived with two different bracketings:
\begin{itemize}
  \item \texttt{(1 + 2) * 3}, which would evaluate to 9.
  \item \texttt{1 + (2 * 3)}, which would evaluate to 7.
\end{itemize}

To disambiguate these cases, \ASLvOneterm{} relies on the concept of an LR(1)
grammar~\cite{ASU86}.
In brief, an LR(1) grammar can be considered both a subcategory of context-free
grammars  and a specific kind of shift-reduce parsing algorithm.
Use of an LR(1) grammar ensures that the grammar of \ASLvOneterm{}:
\begin{itemize}
  \item can be parsed deterministically in a single pass
        (that is, with no backtracking).
  \item can be concisely, precisely, and fully specified in both a
        machine-readable and formal way.
  \item permits use of a variety of parser generators to automatically produce
        a grammar-compliant parser from the machine-readable specification
        of the grammar.
  \item Follows common practices in the parsing of programming languages.
\end{itemize}

We emphasize that the use of an LR(1) parsing algorithm never affects the
validity of a given input text. That is, an LR(1) parsing algorithm will
accept a given input text if and only if it can be derived from the grammar,
and will reject it otherwise.

\section*{Resolving ambiguities: conflicts, token priorities, and associativities}

To precisely define the correct resolution of ambiguities in the \ASLvOneterm{} grammar
(such as the \texttt{1 + 2 * 3} example above), we rely on the shift-reduce parsing
algorithm associated with LR(1) grammars. This is a left-to-right scan of the input
text that keeps track of the stack of already-parsed text, and looks only at the
following token. Based on this token, it chooses to:
\begin{description}
  \item[Shift:] add the token to the stack.
  \item[Reduce:] reorganise the top of the stack (the most recently read elements)
according to a rule of the grammar.
\end{description}

\hypertarget{def-shiftreduceconflict}{}
In cases where the grammar seems ambiguous, it is not clear whether to shift or
reduce. This is called a ``\emph{\shiftreduceconflictterm}''. For the \texttt{1 + 2 * 3} example,
on reading \texttt{*}, it is not clear whether to:
\begin{description}
  \item[Shift:] this would favour the bracketing $\texttt{1 + (2 * 3)}$, as it
  consumes more of the input before performing a bracketing.
  \item[Reduce:] this would favour the bracketing $\texttt{(1 + 2) * 3}$, as it
  reduces (that is, brackets) the existing stack of $\texttt{expr \Tplus{} expr}$ to just
  $\texttt{expr}$ .
\end{description}

To resolve this, we need a relative priority of the existing stack
($\texttt{expr \Tplus{} expr}$) vs. the next token ($\Tmul$).
We use a standard approach here:
\begin{itemize}
\item Assign the existing stack the priority of its rightmost ``terminal''
      (that is, non-symbolic token, in this case $\Tplus$).
\item Assign the next token a priority.
\item Resolve the \shiftreduceconflictterm{} based on priorities:
  \begin{itemize}
  \item If the next token has higher priority, then shift (this is the case for this
        example in ASL, so multiplication is bracketed more tightly than addition).
  \item If the next token has lower priority, then reduce.
  \item Otherwise, the two tokens have the same priority, and the \shiftreduceconflictterm{} is resolved
        by assigning each priority an associativity:
        \begin{itemize}
        \item If the priority of the two tokens has left-associativity, reduce
        (favouring the left-associated bracketing).
        \item If the priority of the two tokens has right-associativity, shift
        (favouring the right-associated bracketing).
        \item Otherwise, the tokens are non-associative, and a parsing error is produced.
        \end{itemize}
  \end{itemize}
\end{itemize}

\ChapterOutline
\begin{itemize}
  \item \secref{InlinedDerivations} defines \emph{inlined derivations};
  \item \secref{ParametricProductions} defines \emph{parametric productions};
  \item \secref{ASLParametricProductions} introduces the parametric productions
        used for the ASL grammar;
  \item \secref{ASLGrammar} defines the ASL grammar in terms of derivations rules;
  \item \secref{ParseTrees} defines parse trees;
  \item \secref{PriorityAndAssociativity} defines the priority and associativity
        of tokens; and
  \item \secref{ParsingExamples} presents examples showing how the grammar and
        the priority and associativity rules are used to parse ASL text.
\end{itemize}

\section{Inlined Derivations\label{sec:InlinedDerivations}}
Context-free grammars consist of a list of \emph{derivations} $N \derives \KleeneStar{S}$
where $N$ is a non-terminal symbol and $S$ is a list of non-terminal symbols and terminal symbols,
which correspond to tokens.
We refer to a list of such symbols as a \emph{sentence}.
\hypertarget{def-emptysentence}{}
A special form of a sentence is the \emph{empty sentence}, written $\emptysentence$.

As commonly done, we aggregate all derivations associated with the same non-terminal symbol
by writing $N \derives R_1 \;|\; \ldots \;|\; R_k$.
We refer to the right-hand-side sentences $R_{1..k}$ as the \emph{alternatives} of $N$.

Our grammar contains another form of derivation --- \emph{inlined derivation} ---
written as $N \derivesinline R_1 \;|\; \ldots \;|\; R_k$.
Expanding an inlined derivation consists of replacing each instance of $N$
in a right-hand-side sentence of a derivation with each of $R_{1..k}$, thereby
creating $k$ variations of it (and removing $N \derivesinline R_1 \;|\; \ldots \;|\; R_k$
from the set of derivations).

For example, consider the derivation
\begin{flalign*}
\Nexpr \derives\ & \Nexpr \parsesep \Nbinop \parsesep \Nexpr &
\end{flalign*}
coupled with the derivation
\begin{flalign*}
\Nbinop \derives\ & \Tand \;|\; \Tband \;|\; \Tbor \;|\; \Tbeq \;|\; \Tdiv \;|\; \Tdivrm \;|\; \Txor \;|\; \Teqop \;|\; \Tneq &\\
                      |\ & \Tgt \;|\; \Tgeq \;|\; \Tbimpl \;|\; \Tlt \;|\; \Tleq \;|\; \Tplus \;|\; \Tminus \;|\; \Tmod \;|\; \Tmul &\\
                      |\ & \Tor \;|\; \Trdiv \;|\; \Tshl \;|\; \Tshr \;|\; \Tpow \;|\; \Tcoloncolon \;|\; \Tplusplus
\end{flalign*}

A grammar containing these two derivations results in \shiftreduceconflictsterm{}.
Resolving these \shiftreduceconflictsterm{} is done by associating priority levels to each of the binary operators
and creating a version of the first derivation for each binary operator:
\begin{flalign*}
\Nexpr \derives\ & \Nexpr \parsesep \Tand \parsesep \Nexpr & \\
              |\ & \Nexpr \parsesep \Tband \parsesep \Nexpr & \\
              |\ & \Nexpr \parsesep \Tbor \parsesep \Nexpr & \\
              \ldots \\
              |\ & \Nexpr \parsesep \Tplusplus \parsesep \Nexpr &
\end{flalign*}

By defining the derivations of $\Nbinop$ as inlined, we achieve the same effect more compactly:
\begin{flalign*}
\Nbinop \derivesinline\ & \Tand \;|\; \Tband \;|\; \Tbor \;|\; \Tbeq \;|\; \Tdiv \;|\; \Tdivrm \;|\; \Txor \;|\; \Teqop \;|\; \Tneq &\\
                      |\ & \Tgt \;|\; \Tgeq \;|\; \Tbimpl \;|\; \Tlt \;|\; \Tleq \;|\; \Tplus \;|\; \Tminus \;|\; \Tmod \;|\; \Tmul &\\
                      |\ & \Tor \;|\; \Trdiv \;|\; \Tshl \;|\; \Tshr \;|\; \Tpow \;|\; \Tcoloncolon \;|\; \Tplusplus
\end{flalign*}

Barring mutually-recursive derivations involving inlined derivations, it is possible to expand
all inlined derivations to obtain a context-free grammar without any inlined derivations.

\section{Parametric Productions\label{sec:ParametricProductions}}
A parametric production has the form
$N(p_{1..m}) \derives R_1 \;|\; \ldots \;|\; R_k$
where $p_{1..m}$ are place holders for grammar symbols and may appear in any of the alternatives $R_{1..k}$.
We refer to $N(p_{1..m})$ as a \emph{parametric non-terminal}.

Given sentences $S_{1..m}$, we can expand $N(p_{1..m}) \derives R_1 \;|\; \ldots \;|\; R_k$
by creating a unique symbol for $N(S_{1..m})$, denoted as $\uniquesymb{N(S_{1..m})}$, defining the
derivations
\[
  \uniquesymb{N(S_{1..m})} \derives R_1[S_1/p_1,\ldots,S_m/p_m] \;|\; \ldots \;|\; R_k[S_1/p_1,\ldots,S_m/p_m]
\]
where for each $i= 1..k$, $R_i[S_1/p_1,\ldots,S_m/p_m]$ means replacing each instance of $p_j$ with $S_j$, for each $j=1..m$.
Then, each instance of $N(S_{1..m})$ in the grammar is replaced by $\uniquesymb{N(S_{1..m})}$.
Expanding all instances of parametric productions results in a grammar without any parametric productions.

We note that a parametric production can be either a normal derivation or an inlined derivation.

\ExampleDef{Parametric List of Global Declaration}
The derivation for a list of ASL global declarations is as follows:
\begin{flalign*}
\Nspec \derives\ & \maybeemptylist{\Ndecl} &
\end{flalign*}
It is defined via the parametric production for possibly-empty lists:
\begin{flalign*}
\maybeemptylist{x}   \derives\ & \emptysentence \;|\; x \parsesep \maybeemptylist{x} &\\
\end{flalign*}

Expanding $\maybeemptylist{\Ndecl}$ produces the following derivations for a new unique symbol.
That is, a symbol that does not appear anywhere else in the grammar.
In this example we will choose $\uniquesymb{\maybeemptylist{\Ndecl}}$ to be the symbol $\Ndecllist$.
The result of the expansion is then:
\begin{flalign*}
\Ndecllist   \derives\ & \emptysentence \;|\; \Ndecl \parsesep \Ndecllist &\\
\end{flalign*}
The new symbol is substituted anywhere $\maybeemptylist{\Ndecl}$ appears in the original grammar,
which results in the following derivation replacing the original derivation for $\Nspec$:
\begin{flalign*}
\Nspec \derives\ & \Ndecllist &
\end{flalign*}

\section{ASL Parametric Productions\label{sec:ASLParametricProductions}}
We define the following parametric productions for various types of lists and optional productions.

\paragraph{Optional Symbol}
\hypertarget{def-option}{}
\begin{flalign*}
\option{x}   \derives\ & \emptysentence \;|\; x &\\
\end{flalign*}

\paragraph{Possibly-empty List}
\hypertarget{def-maybeemptylist}{}
\begin{flalign*}
\maybeemptylist{x}   \derives\ & \emptysentence \;|\; x \parsesep \maybeemptylist{x} &\\
\end{flalign*}

\paragraph{Non-empty List}
\hypertarget{def-listone}{}
\begin{flalign*}
\ListOne{x}   \derives\ & x \;|\; x \parsesep \ListOne{x}&\\
\end{flalign*}

\paragraph{Non-empty Comma-separated List}
\hypertarget{def-clistone}{}
\begin{flalign*}
\ClistOne{x}   \derives\ & x \;|\; x \parsesep \Tcomma \parsesep \ClistOne{x} &\\
\end{flalign*}

\paragraph{Possibly-empty Comma-separated List}
\hypertarget{def-clistzero}{}
\begin{flalign*}
\ClistZero{x}   \derives \ & \emptysentence \;|\; \ClistOne{x} &\\
\end{flalign*}

\paragraph{Comma-separated List With At Least Two Elements}
\hypertarget{def-clisttwo}{}
\begin{flalign*}
\Clisttwo{x}   \derives \ & x \parsesep \Tcomma \parsesep \ClistOne{x} &\\
\end{flalign*}

\paragraph{Possibly-empty Parenthesized, Comma-separated List}
\hypertarget{def-plistzero}{}
\begin{flalign*}
\PlistZero{x}   \derivesinline \ & \Tlpar \parsesep \ClistZero{x} \parsesep \Trpar &\\
\end{flalign*}

\paragraph{Parenthesized Comma-separated List With At Least Two Elements}
\hypertarget{def-plisttwo}{}
\begin{flalign*}
\Plisttwo{x}   \derivesinline \ & \Tlpar \parsesep \Clisttwo{x} \parsesep \Trpar &\\
\end{flalign*}

\paragraph{Non-empty Comma-separated Trailing List}
\hypertarget{def-tclistone}{}
\begin{flalign*}
\TClistOne{x}   \derives\ & x \parsesep \option{\Tcomma} &\\
                          |\  & x \parsesep \Tcomma \parsesep \TClistOne{x}
\end{flalign*}

\paragraph{Possibly-empty Comma-separated Trailing List}
\hypertarget{def-tclistzero}{}
\begin{flalign*}
\TClistZero{x}   \derives \ & \option{\TClistOne{x}} &\\
\end{flalign*}

\section{ASL Grammar\label{sec:ASLGrammar}}
We now present the list of derivations for the ASL Grammar where the start non-terminal is $\Nspec$.
%
The derivations allow certain parse trees where lists may have invalid sizes.
Those parse trees must be rejected in a later phase.

\hypertarget{def-tunops}{}
Notice that two of the derivations (for $\Nexprpattern$ and for $\Nexpr$) end with
$\priority{\Tunops}$.
This is a priority annotation, which is not part of the right-hand-side sentence, and is explained in \secref{PriorityAndAssociativity}
and can be ignored upon first reading.

For brevity, tokens are presented via their label only, dropping their associated value.
For example, instead of $\Tidentifier(\id)$, we simply write $\Tidentifier$.

\hypertarget{def-nspec}{}
\begin{flalign*}
\Nspec   \derives\ & \maybeemptylist{\Ndecl} &
\end{flalign*}

\hypertarget{def-ndecl}{}
\begin{flalign*}
\Ndecl  \derives \ & \Npuritykeyword \parsesep \Noverride \parsesep \Tfunc \parsesep \Tidentifier \parsesep \Nparamsopt \parsesep \Nfuncargs &\\
    &\wrappedline\ \Nreturntype \parsesep \Nrecurselimit \parsesep \Nfuncbody\\
|\ & \Nqualifier \parsesep \Noverride \parsesep \Tfunc \parsesep \Tidentifier \parsesep \Nparamsopt \parsesep \Nfuncargs &\\
    &\wrappedline\ \Nrecurselimit \parsesep \Nfuncbody &\\
|\ & \Noverride \parsesep \Taccessor \parsesep \Tidentifier \parsesep \Nparamsopt \parsesep \Nfuncargs \parsesep \Tbeq \parsesep \Tidentifier \parsesep \Nasty &\\
   & \wrappedline\ \Naccessorbody &\\
|\ & \Ttype \parsesep \Tidentifier \parsesep \Tof \parsesep \Ntydecl \parsesep \Nsubtypeopt \parsesep \Tsemicolon &\\
|\ & \Ttype \parsesep \Tidentifier \parsesep \Nsubtype \parsesep \Tsemicolon &\\
|\ & \Nglobaldeclkeywordnonconfig \parsesep \Tidentifier \parsesep \option{\Tcolon \parsesep \Nty} &\\
   & \wrappedline\ \Teq \parsesep \Nexpr \parsesep \Tsemicolon &\\
|\ & \Tconfig \parsesep \Tidentifier \parsesep \Tcolon \parsesep \Nty \parsesep \Teq \parsesep \Nexpr \parsesep \Tsemicolon &\\
|\ & \Tvar \parsesep \Tidentifier \parsesep \Tcolon \parsesep \Ntyorcollection \parsesep \Tsemicolon &\\
|\ & \Tvar \parsesep \Clisttwo{\Tidentifier} \parsesep \Nasty \parsesep \Tsemicolon &\\
|\ & \Nglobaldeclkeyword \parsesep \Tidentifier \parsesep \Tcolon \parsesep \Nty &\\
   & \wrappedline\ \Teq \parsesep \Nelidedparamcall \parsesep \Tsemicolon &\\
|\ & \Tpragma \parsesep \Tidentifier \parsesep \ClistZero{\Nexpr} \parsesep \Tsemicolon&
\end{flalign*}

\hypertarget{def-nrecurselimit}{}
\begin{flalign*}
\Nrecurselimit   \derives \ & \Trecurselimit \parsesep \Nexpr &\\
|\                          & \emptysentence &\\
\end{flalign*}

\hypertarget{def-nsubtype}{}
\begin{flalign*}
\Nsubtype \derives\ & \Tsubtypes \parsesep \Tidentifier \parsesep \Twith \parsesep \Nfields &\\
          |\        & \Tsubtypes \parsesep \Tidentifier &\\
\end{flalign*}

\hypertarget{def-nsubtypeopt}{}
\begin{flalign*}
\Nsubtypeopt           \derives \ & \option{\Nsubtype} &
\end{flalign*}

\hypertarget{def-ntypedidentifier}{}
\begin{flalign*}
\Ntypedidentifier \derives \ & \Tidentifier \parsesep \Nasty &
\end{flalign*}

\hypertarget{def-nopttypeidentifier}{}
\begin{flalign*}
\Nopttypedidentifier \derives \ & \Tidentifier \parsesep \option{\Nasty} &
\end{flalign*}

\hypertarget{def-nasty}{}
\begin{flalign*}
\Nasty \derives \ & \Tcolon \parsesep \Nty &
\end{flalign*}

\hypertarget{def-nreturntype}{}
\begin{flalign*}
\Nreturntype        \derives \ & \Tarrow \parsesep \Nty &
\end{flalign*}

\hypertarget{def-nparamsopt}{}
\begin{flalign*}
\Nparamsopt \derives \ & \emptysentence &\\
                    |\ & \Tlbrace \parsesep \ClistZero{\Nopttypedidentifier} \parsesep \Trbrace &
\end{flalign*}

\hypertarget{def-ncall}{}
\begin{flalign*}
\Ncall \derives \
     & \Tidentifier \parsesep \PlistZero{\Nexpr} &\\
  |\ & \Tidentifier \parsesep \Tlbrace \parsesep \ClistOne{\Nexpr} \parsesep \Trbrace &\\
  |\ & \Tidentifier \parsesep \Tlbrace \parsesep \ClistOne{\Nexpr} \parsesep \Trbrace \parsesep \PlistZero{\Nexpr} &
\end{flalign*}

\hypertarget{def-nelidedparamcall}{}
\begin{flalign*}
\Nelidedparamcall \derives \
     & \Tidentifier \parsesep \Tlbrace \parsesep \Trbrace &\\
  |\ & \Tidentifier \parsesep \Tlbrace \parsesep \Trbrace \parsesep \PlistZero{\Nexpr} &\\
  |\ & \Tidentifier \parsesep \Tlbrace \parsesep \Tcomma \parsesep \ClistOne{\Nexpr} \parsesep \Trbrace &\\
  |\ & \Tidentifier \parsesep \Tlbrace \parsesep \Tcomma \parsesep \ClistOne{\Nexpr} \parsesep \Trbrace \parsesep \PlistZero{\Nexpr}&
\end{flalign*}

\hypertarget{def-nfuncargs}{}
\begin{flalign*}
    \Nfuncargs          \derives \ & \PlistZero{\Ntypedidentifier} &
\end{flalign*}

\hypertarget{def-nfuncbody}{}
\begin{flalign*}
\Nfuncbody          \derives \ & \Tbegin \parsesep \Nstmtlist \parsesep \Tend \parsesep \Tsemicolon &
\end{flalign*}

\hypertarget{def-nignoredoridentifier}{}
\begin{flalign*}
\Nignoredoridentifier \derives \ & \Tminus \;|\; \Tidentifier &
\end{flalign*}

\hypertarget{def-naccessorbody}{}
\begin{flalign*}
\Naccessorbody \derives \ & \Tbegin \parsesep \Naccessors \parsesep \Tend \parsesep \Tsemicolon&
\end{flalign*}

\hypertarget{def-naccessors}{}
\begin{flalign*}
   \Naccessors \derives \ &
      \Nisreadonly \parsesep \Tgetter \parsesep \Nstmtlist \parsesep \Tend \parsesep \Tsemicolon \parsesep \\ & \wrappedline
      \Tsetter \parsesep \Nstmtlist \parsesep \Tend \parsesep \Tsemicolon &\\
   |\ & \Tsetter \parsesep \Nstmtlist \parsesep \Tend \parsesep \Tsemicolon \parsesep \\ & \wrappedline
        \Nisreadonly \parsesep \Tgetter \parsesep \Nstmtlist \parsesep \Tend \parsesep \Tsemicolon &\\
\end{flalign*}

\hypertarget{def-nqualifier}{}
\begin{flalign*}
\Nqualifier \derivesinline\ & \emptysentence \;|\; \Tpure \;|\;\Treadonly \;|\; \Tnoreturn &
\end{flalign*}

\hypertarget{def-npuritykeyword}{}
\begin{flalign*}
\Npuritykeyword \derivesinline\ & \emptysentence \;|\; \Tpure \;|\;\Treadonly &
\end{flalign*}

\hypertarget{def-nisreadonly}{}
\begin{flalign*}
\Nisreadonly \derivesinline\ & \emptysentence \;|\; \Treadonly &
\end{flalign*}

\hypertarget{def-noverride}{}
\begin{flalign*}
\Noverride \derivesinline\ & \emptysentence \;|\; \Timpdef \;|\;\Timplementation &
\end{flalign*}

\vspace*{-\baselineskip}
\hypertarget{def-nlocaldeclkeyword}{}
\begin{flalign*}
\Nlocaldeclkeyword \derives \ & \Tlet \;|\; \Tconstant \;|\; \Tvar &
\end{flalign*}

\hypertarget{def-nglobaldeclkeywordnonconfig}{}
\begin{flalign*}
\Nglobaldeclkeywordnonconfig \derives \ & \Tlet \;|\; \Tconstant \;|\; \Tvar &
\end{flalign*}

\hypertarget{def-nglobaldeclkeyword}{}
\begin{flalign*}
\Nglobaldeclkeyword \derives \ & \Nglobaldeclkeywordnonconfig \;|\; \Tconfig &
\end{flalign*}

\hypertarget{def-ndirection}{}
\begin{flalign*}
\Ndirection \derives \ & \Tto \;|\; \Tdownto &
\end{flalign*}

\hypertarget{def-ncasealtlist}{}
\begin{flalign*}
\Ncasealtlist \derives \ & \ClistOne{\Ncasealt} &\\
\end{flalign*}

\hypertarget{def-ncasealt}{}
\begin{flalign*}
\Ncasealt \derives \ & \Twhen \parsesep \Npatternlist \parsesep \option{\Twhere \parsesep \Nexpr} \parsesep \Tarrow \parsesep \Nstmtlist &
\end{flalign*}

\hypertarget{def-notherwiseopt}{}
\begin{flalign*}
\Notherwiseopt \derives\ & \Totherwise \parsesep \Tarrow \parsesep \Nstmtlist &\\
                      |\ & \emptysentence &\\
\end{flalign*}

\hypertarget{def-ncatcher}{}
\begin{flalign*}
\Ncatcher \derives      \ & \Twhen \parsesep \Tidentifier \parsesep \Tcolon \parsesep \Nty \parsesep \Tarrow \parsesep \Nstmtlist &\\
          |\              & \Twhen \parsesep \Nty \parsesep \Tarrow \parsesep \Nstmtlist &\\
\end{flalign*}

\hypertarget{def-nlooplimit}{}
\begin{flalign*}
\Nlooplimit \derives    \ & \Tlooplimit \parsesep \Nexpr &\\
          |\              & \emptysentence &\\
\end{flalign*}

\hypertarget{def-nstmt}{}
\begin{flalign*}
\Nstmt \derives \ & \Tif \parsesep \Nexpr \parsesep \Tthen \parsesep \Nstmtlist \parsesep \Nselse \parsesep \Tend \parsesep \Tsemicolon &\\
|\ & \Tcase \parsesep \Nexpr \parsesep \Tof \parsesep \Ncasealtlist \parsesep \Tend \parsesep \Tsemicolon &\\
|\ & \Tcase \parsesep \Nexpr \parsesep \Tof \parsesep \Ncasealtlist \parsesep \Totherwise \parsesep \Tarrow &\\
   & \wrappedline\ \Nstmtlist \parsesep \Tend \parsesep \Tsemicolon &\\
|\ & \Twhile \parsesep \Nexpr \parsesep \Nlooplimit \parsesep \Tdo \parsesep \Nstmtlist \parsesep \Tend \parsesep \Tsemicolon &\\
|\ & \Tfor \parsesep \Tidentifier \parsesep \Teq \parsesep \Nexpr \parsesep \Ndirection \parsesep
                    \Nexpr \parsesep \Nlooplimit \parsesep \Tdo &\\
                    & \wrappedline\ \Nstmtlist \parsesep \Tend \parsesep \Tsemicolon &\\
|\ & \Ttry \parsesep \Nstmtlist \parsesep \Tcatch \parsesep \ListOne{\Ncatcher} \parsesep \Notherwiseopt \parsesep \Tend \parsesep \Tsemicolon &\\
|\ & \Tpass \parsesep \Tsemicolon &\\
|\ & \Treturn \parsesep \option{\Nexpr} \parsesep \Tsemicolon &\\
|\ & \Ncall \parsesep \Tsemicolon &\\
|\ & \Tassert \parsesep \Nexpr \parsesep \Tsemicolon &\\
|\ & \Nlocaldeclkeyword \parsesep \Ndeclitem \parsesep \option{\Nasty} \parsesep \Teq \parsesep \Nexpr \parsesep \Tsemicolon &\\
|\ & \Nlexpr \parsesep \Teq \parsesep \Nexpr \parsesep \Tsemicolon &\\
|\ & \Ncall \parsesep \Nsetteraccess \parsesep \Teq \parsesep \Nexpr \parsesep \Tsemicolon &\\
|\ & \Ncall \parsesep \Nsetteraccess \parsesep \Nslices \parsesep \Teq \parsesep \Nexpr \parsesep \Tsemicolon &\\
|\ & \Ncall \parsesep \Tdot \parsesep \Tlbracket \parsesep \Clisttwo{{\Tidentifier}} \parsesep \Trbracket \parsesep \Teq \parsesep \Nexpr \parsesep \Tsemicolon &\\
|\ & \Nlocaldeclkeyword \parsesep \Ndeclitem \parsesep \Nasty \parsesep \Teq \parsesep \Nelidedparamcall \parsesep \Tsemicolon &\\
|\ & \Tvar \parsesep \Ndeclitem \parsesep \Nasty \parsesep \Tsemicolon &\\
|\ & \Tvar \parsesep \Clisttwo{\Tidentifier} \parsesep \Nasty \parsesep \Tsemicolon &\\
|\ & \Tprint \parsesep \ClistZero{\Nexpr} \parsesep \Tsemicolon &\\
|\ & \Tprintln \parsesep \ClistZero{\Nexpr} \parsesep \Tsemicolon &\\
|\ & \Tunreachable \parsesep \Tsemicolon &\\
|\ & \Trepeat \parsesep \Nstmtlist \parsesep \Tuntil \parsesep \Nexpr \parsesep \Nlooplimit \parsesep \Tsemicolon &\\
|\ & \Tthrow \parsesep \Nexpr \parsesep \Tsemicolon &\\
|\ & \Tpragma \parsesep \Tidentifier \parsesep \ClistZero{\Nexpr} \parsesep \Tsemicolon &
\end{flalign*}

\hypertarget{def-nstmtlist}{}
\begin{flalign*}
\Nstmtlist \derives \ & \ListOne{\Nstmt} &
\end{flalign*}

\hypertarget{def-nselse}{}
\begin{flalign*}
\Nselse \derives\ & \Telseif \parsesep \Nexpr \parsesep \Tthen \parsesep \Nstmtlist \parsesep \Nselse &\\
|\ & \Telse \parsesep \Nstmtlist &\\
|\ & \emptysentence &
\end{flalign*}

\hypertarget{def-nlexpr}{}
\begin{flalign*}
\Nlexpr \derives\
   & \Tminus &\\
|\ & \Nbasiclexpr &\\
|\ & \Tlpar \parsesep \Clisttwo{\Ndiscardorbasiclexpr} \parsesep \Trpar &\\
|\ & \Tidentifier \parsesep \Tdot \parsesep \Tlbracket \parsesep \Clisttwo{\Tidentifier} \parsesep \Trbracket &\\
|\ & \Tidentifier \parsesep \Tdot \parsesep \Tlpar \parsesep \Clisttwo{\Ndiscardoridentifier} \parsesep \Trpar &
\end{flalign*}

\hypertarget{def-nbasiclexpr}{}
\begin{flalign*}
\Nbasiclexpr \derives\
   & \Tidentifier \parsesep \Naccess &\\
|\ & \Tidentifier \parsesep \Naccess \parsesep \Nslices &
\end{flalign*}

\hypertarget{def-naccess}{}
\begin{flalign*}
\Naccess \derives\
   & \emptysentence  &\\
|\ & \Tdot \parsesep \Tidentifier \parsesep \Naccess &\\
|\ & \Tllbracket \parsesep \Nexpr \parsesep \Trrbracket \parsesep \Naccess &
\end{flalign*}

\hypertarget{def-ndiscardorbasiclexpr}{}
\begin{flalign*}
\Ndiscardorbasiclexpr \derives\ & \Tminus \;|\; \Nbasiclexpr &
\end{flalign*}

\hypertarget{def-ndiscardoridentifier}{}
\begin{flalign*}
\Ndiscardoridentifier \derives \ & \Tminus \;|\; \Tidentifier &
\end{flalign*}

\hypertarget{def-nsetteraccess}{}
\begin{flalign*}
\Nsetteraccess \derives \
   & \emptysentence &\\
|\ & \Tdot \parsesep \Tidentifier \parsesep \Nsetteraccess &
\end{flalign*}

A $\Ndeclitem$ is another kind of left-hand-side expression,
which appears only in declarations. It cannot have setter calls or set record fields,
it must declare a new variable.
\hypertarget{def-ndeclitem}{}
\begin{flalign*}
\Ndeclitem \derives\
   & \Tidentifier &\\
|\ & \Plisttwo{\Nignoredoridentifier}  &
\end{flalign*}

\hypertarget{def-nintconstraintsopt}{}
\begin{flalign*}
\Nconstraintkindopt \derives \ & \Nconstraintkind \;|\; \emptysentence &
\end{flalign*}

\hypertarget{def-nintconstraints}{}
\begin{flalign*}
\Nconstraintkind \derives \ &
       \Tlbrace \parsesep \ClistOne{\Nintconstraint} \parsesep \Trbrace &\\
  |\ & \Tlbrace \parsesep \Trbrace &
\end{flalign*}

\hypertarget{def-nintconstraint}{}
\begin{flalign*}
\Nintconstraint \derives \ & \Nexpr &\\
|\ & \Nexpr \parsesep \Tslicing \parsesep \Nexpr &
\end{flalign*}

Pattern expressions ($\Nexprpattern$), defined by the following derivations, are similar to regular expressions  ($\Nexpr$),
except they do not derive tuples, which are the last derivation for $\Nexpr$.

\hypertarget{def-nexprpattern}{}
\begin{flalign*}
\Nexprpattern \derives\ & \Nvalue &\\
                    |\  & \Tidentifier &\\
                    |\  & \Nexprpattern \parsesep \Nbinop \parsesep \Nexpr &\\
                    |\  & \Nunop \parsesep \Nexpr & \priority{\Tunops}\\
                    |\  & \Tif \parsesep \Nexpr \parsesep \Tthen \parsesep \Nexpr \parsesep \Telse \parsesep \Nexpr &\\
                    |\  & \Ncall &\\
                    |\  & \Nexprpattern \parsesep \Nslices &\\
                    |\  & \Nexprpattern \parsesep \Tllbracket \parsesep \Nexpr \parsesep \Trrbracket &\\
                    |\  & \Nexprpattern \parsesep \Tdot \parsesep \Tidentifier&\\
                    |\  & \Nexprpattern \parsesep \Tdot \parsesep \Tlbracket \parsesep \ClistOne{\Tidentifier} \parsesep \Trbracket &\\
                    |\  & \Nexprpattern \parsesep \Tas \parsesep \Nty &\\
                    |\  & \Nexprpattern \parsesep \Tas \parsesep \Nconstraintkind &\\
                    |\  & \Nexpr \parsesep \Tin \parsesep \Npatternset &\\
                    |\  & \Nexpr \parsesep \Teqop \parsesep \Tmasklit &\\
                    |\  & \Nexpr \parsesep \Tneq \parsesep \Tmasklit &\\
                    |\  & \Tarbitrary \parsesep \Tcolon \parsesep \Nty &\\
                    |\  & \Tidentifier \parsesep \Tlbrace \parsesep \Tminus \parsesep \Trbrace & \mathhypertarget{any-pattern}\\
                    |\  & \Tidentifier \parsesep \Tlbrace \parsesep \ClistOne{\Nfieldassign} \parsesep \Trbrace &\\
                    |\  & \Tlpar \parsesep \Nexprpattern \parsesep \Trpar &
\end{flalign*}

\hypertarget{def-npatternset}{}
\begin{flalign*}
\Npatternset \derives \  & \Tbnot \parsesep \Tlbrace \parsesep \Npatternlist \parsesep \Trbrace &\\
                  |\    & \Tlbrace \parsesep \Npatternlist \parsesep \Trbrace &
\end{flalign*}

\hypertarget{def-npatternlist}{}
\begin{flalign*}
\Npatternlist \derives \ & \ClistOne{\Npattern} &
\end{flalign*}

\hypertarget{def-npattern}{}
\begin{flalign*}
\Npattern \derives\ & \Nexprpattern &\\
                |\  & \Nexprpattern \parsesep \Tslicing \parsesep \Nexpr &\\
                |\  & \Tminus &\\
                |\  & \Tleq \parsesep \Nexpr &\\
                |\  & \Tgeq \parsesep \Nexpr &\\
                |\  & \Tmasklit &\\
                |\  & \Plisttwo{\Npattern} &\\
                |\  & \Npatternset &
\end{flalign*}

\hypertarget{def-nfields}{}
\begin{flalign*}
\Nfields \derives \ & \Tlbrace \parsesep \Tminus \parsesep \Trbrace &\\
                  |\ & \Tlbrace \parsesep \TClistOne{\Ntypedidentifier} \parsesep \Trbrace &
\end{flalign*}

\hypertarget{def-nslices}{}
\begin{flalign*}
\Nslices \derives \ & \Tlbracket \parsesep \ClistOne{\Nslice} \parsesep \Trbracket &
\end{flalign*}

\hypertarget{def-nslice}{}
\begin{flalign*}
\Nslice \derives \ & \Nexpr &\\
              |\  & \Nexpr \parsesep \Tcolon \parsesep \Nexpr &\\
              |\  & \Nexpr \parsesep \Tpluscolon \parsesep \Nexpr &\\
              |\  & \Nexpr \parsesep \Tstarcolon \parsesep \Nexpr &\\
              |\  & \Tcolon \parsesep \Nexpr &
\end{flalign*}

\hypertarget{def-nbitfields}{}
\begin{flalign*}
\Nbitfields \derives \ & \Tlbrace \parsesep \TClistZero{\Nbitfield} \parsesep \Trbrace &
\end{flalign*}

\hypertarget{def-nbitfield}{}
\begin{flalign*}
\Nbitfield \derives \ & \Nslices \parsesep \Tidentifier &\\
                  |\ & \Nslices \parsesep \Tidentifier \parsesep \Nbitfields &\\
                  |\ & \Nslices \parsesep \Tidentifier \parsesep \Tcolon \parsesep \Nty &
\end{flalign*}

\hypertarget{def-nty}{}
\begin{flalign*}
\Nty \derives\ & \Tinteger \parsesep \Nconstraintkindopt &\\
            |\ & \Treal &\\
            |\ & \Tstring &\\
            |\ & \Tboolean &\\
            |\ & \Tbit &\\
            |\ & \Tbits \parsesep \Tlpar \parsesep \Nexpr \parsesep \Trpar \parsesep \option{\Nbitfields} &\\
            |\ & \Tlpar \parsesep \Nty \parsesep \Trpar &\\
            |\ & \Plisttwo{\Nty} &\\
            |\ & \Tidentifier &\\
            |\ & \Tarray \parsesep \Tllbracket \parsesep \Nexpr \parsesep \Trrbracket \parsesep \Tof \parsesep \Nty &
\end{flalign*}

\hypertarget{def-ntydecl}{}
\begin{flalign*}
\Ntydecl \derives\ & \Nty &\\
            |\ & \Tenumeration \parsesep \Tlbrace \parsesep \TClistOne{\Tidentifier} \parsesep \Trbrace &\\
            |\ & \Trecord \parsesep \Nfields &\\
            |\ & \Texception \parsesep \Nfields &
\end{flalign*}

\hypertarget{def-nfieldassign}{}
\begin{flalign*}
\Nfieldassign \derives \ & \Tidentifier \parsesep \Teq \parsesep \Nexpr &
\end{flalign*}

\hypertarget{def-ntyorcollection}{}
\begin{flalign*}
  \Ntyorcollection \derives \ & \Nty \;|\; \Tcollection \parsesep \Nfields &\\
\end{flalign*}

\hypertarget{def-nexpr}{}
\begin{flalign*}
\Nexpr \derives\  & \Nvalue &\\
                    |\  & \Tidentifier &\\
                    |\  & \Nexpr \parsesep \Nbinop \parsesep \Nexpr &\\
                    |\  & \Nunop \parsesep \Nexpr & \priority{\Tunops}\\
                    |\  & \Tif \parsesep \Nexpr \parsesep \Tthen \parsesep \Nexpr \parsesep \Telse \parsesep \Nexpr &\\
                    |\  & \Ncall &\\
                    |\  & \Nexpr \parsesep \Nslices &\\
                    |\  & \Nexpr \parsesep \Tllbracket \parsesep \Nexpr \parsesep \Trrbracket &\\
                    |\  & \Nexpr \parsesep \Tdot \parsesep \Tidentifier&\\
                    |\  & \Nexpr \parsesep \Tdot \parsesep \Tlbracket \parsesep \ClistOne{\Tidentifier} \parsesep \Trbracket &\\
                    |\  & \Nexpr \parsesep \Tas \parsesep \Nty &\\
                    |\  & \Nexpr \parsesep \Tas \parsesep \Nconstraintkind &\\
                    |\  & \Nexpr \parsesep \Tin \parsesep \Npatternset &\\
                    |\  & \Nexpr \parsesep \Teqop \parsesep \Tmasklit &\\
                    |\  & \Nexpr \parsesep \Tneq \parsesep \Tmasklit &\\
                    |\  & \Tarbitrary \parsesep \Tcolon \parsesep \Nty &\\
                    |\  & \Tidentifier \parsesep \Tlbrace \parsesep \Tminus \parsesep \Trbrace &\\
                    |\  & \Tidentifier \parsesep \Tlbrace \parsesep \ClistOne{\Nfieldassign} \parsesep \Trbrace &\\
                    |\  & \Tlpar \parsesep \Nexpr \parsesep \Trpar &\\
                    |\  & \Plisttwo{\Nexpr} &
\end{flalign*}

\hypertarget{def-nvalue}{}
\begin{flalign*}
\Nvalue \derives      \ & \Tintlit &\\
                     |\ & \Tboollit &\\
                     |\ & \Treallit &\\
                     |\ & \Tbitvectorlit &\\
                     |\ & \Tstringlit &
\end{flalign*}

\hypertarget{def-nunop}{}
\begin{flalign*}
\Nunop \derivesinline\ & \Tbnot \;|\; \Tminus \;|\; \Tnot &
\end{flalign*}

\hypertarget{def-nbinop}{}
\begin{flalign*}
\Nbinop \derivesinline\ & \Tand \;|\; \Tband \;|\; \Tbor \;|\; \Tbeq \;|\; \Tdiv \;|\; \Tdivrm \;|\; \Txor \;|\; \Teqop \;|\; \Tneq &\\
                     |\ & \Tgt \;|\; \Tgeq \;|\; \Tbimpl \;|\; \Tlt \;|\; \Tleq \;|\; \Tplus \;|\; \Tminus \;|\; \Tmod \;|\; \Tmul &\\
                     |\ & \Tor \;|\; \Trdiv \;|\; \Tshl \;|\; \Tshr \;|\; \Tpow \;|\; \Tcoloncolon \;|\; \Tplusplus
\end{flalign*}

\section{Parse Trees\label{sec:ParseTrees}}
We now define \emph{parse trees} for the ASL expanded grammar. Those are later used to build Abstract Syntax Trees.

\begin{definition}[Parse Trees]
A \emph{parse tree} has one of the following forms:
\begin{itemize}
  \item A \emph{token node}, given by the token itself, for example, $\Tlexeme(\Tarrow)$ and $\Tidentifier(\id)$;
  \item \hypertarget{def-epsilonnode}{} $\epsilonnode$, which represents the empty sentence --- $\emptysentence$; or
  \item A \emph{non-terminal node} of the form $N(n_{1..k})$ where $N$ is a non-terminal symbol,
        which is said to label the node,
        and $n_{1..k}$ are its children parse nodes,
        for example,
        $\Ndecl(\Tfunc, \Tidentifier(\id), \Nparamsopt, \Nfuncargs, \Nfuncbody)$
        is labeled by $\Ndecl$ and has five children nodes.
\end{itemize}
\end{definition}
(In the literature, parse trees are also referred to as \emph{derivation trees}.)

\begin{definition}[Well-formed Parse Trees]
A parse tree is \emph{well-formed} if its root is labelled by the start non-terminal ($\Nspec$ for ASL)
and each non-terminal node $N(n_{1..k})$ corresponds to a grammar derivation
$N \derives l_{1..k}$ where for each $i \in 1..k$ either:
\begin{itemize}
   \item $n_i$ is a non-terminal node and $l_i$ is its label;
   \item $n_i$ is a token and $l_i$ is equal to $n_i$.
\end{itemize}
A non-terminal node $N(\epsilonnode)$ is well-formed if the grammar includes a derivation
$N \derives \emptysentence$.
\end{definition}

\hypertarget{def-yield}{}
\begin{definition}[Parse Tree Yield]
The \emph{\yield} of a parse tree is the list of tokens
given by an in-order walk of the tree:
\[
\yield(n) \triangleq \begin{cases}
  [t] & n \text{ is a token }t\\
  \emptylist & n = \epsilonnode\\
  \yield(n_1) \concat \ldots \concat \yield(n_k) & n = N(n_{1..k})\\
\end{cases}
\]
\end{definition}

\hypertarget{def-parsenode}{}
We denote the set of well-formed parse trees for a non-terminal symbol $S$ by $\parsenode{S}$.

\hypertarget{def-aslparse}{}
A parser is a function
\[
\aslparse(\KleeneStar{(\Token \setminus \{\Terror\})}) \aslto \parsenode{\Nspec} \cup \{\ParseErrorConfig\}
\]
\hypertarget{def-parseerror}{}
where $\ParseErrorConfig$ stands for a \emph{parse error}.
If $\aslparse(\ts) = n$ then $\yield(n)=\ts$
and if $\aslparse(\ts) = \ParseErrorConfig$ then there is no well-formed tree
$n$ such that $\yield(n)=\ts$.
(Notice that we do not define a parser if $\ts$ is lexically illegal.)

The \emph{language of a grammar} $G$ is defined as follows:
\[
\Lang{G} = \{\yield(n) \;|\; n \text{ is a well-formed parse tree for }G\} \enspace.
\]

\section{Priority and Associativity\label{sec:PriorityAndAssociativity}}
As described in the introduction to \chapref{Syntax}, tokens that may cause a
\shiftreduceconflictterm{} are assigned a priority, and each priority is
assigned an associativity. Tokens that cannot cause \shiftreduceconflictsterm{}
are not assigned priorities

The table below assigns priorities to tokens in descending order, starting from
the highest priority (for $\Tdot$) to the lowest priority (for $\Telse$).
Each priority is assigned an associativity among $\leftassoc$, $\rightassoc$,
and $\nonassoc$.

The rules involving unary operations ($\Nunop$) are not assigned the priority
level of the token corresponding to the unary operation, but rather the priority
level $\Tunops$, as denoted by the notation priority: $\Tunops$ appearing to
their right.
This is a standard way of dealing with unary operations in many programming
languages, which involves defining an artificial token $\Tunops$, which is
never returned by the parser. Note that the $\nonassoc$ associativity of $\Tunops$
does not mean that unary operations cannot be nested. Rather, ASL's grammar cannot
give rise to a \shiftreduceconflictterm{} in which both relevant tokens are
unary operations.
Therefore, no associativity is specified for this impossible case, but for
example \texttt{!!TRUE} still parses as \texttt{!(!TRUE)}.

\begin{center}
\begin{tabular}{lll}
\textbf{Priority level} & \textbf{Tokens} & \textbf{Associativity}\\
\hline
9 & $\Tdot$, $\Tlbracket$, $\Tllbracket$                                          & $\leftassoc$\\
8 & $\Tin$                                                                        & $\nonassoc$\\
7 & $\Tunops$                                                                     & $\nonassoc$\\
6 & $\Tpow$                                                                       & $\leftassoc$\\
5 & $\Tmul$, $\Tdiv$, $\Tdivrm$, $\Trdiv$, $\Tmod$, $\Tshl$, $\Tshr$              & $\leftassoc$\\
4 & $\Tplus$, $\Tminus$, $\Tor$, $\Tand$, $\Txor$, $\Tcoloncolon$, $\Tplusplus$   & $\leftassoc$\\
3 & $\Tgt$, $\Tgeq$, $\Tlt$, $\Tleq$                                              & $\nonassoc$\\
2 & $\Teqop$, $\Tneq$                                                             & $\leftassoc$\\
1 & $\Tbor$, $\Tband$, $\Tbimpl$, $\Tbeq$, $\Tas$                                  & $\leftassoc$\\
0 & $\Telse$                                                                      & $\nonassoc$
\end{tabular}
\end{center}


\section{Parsing Examples\label{sec:ParsingExamples}}

\ExampleDef{Derivation of a Valid Expression}
The expression \texttt{1 + 3}
is transformed by the lexical analysis into:
\[
\Tintlit(1) \parsesep \Tplus \parsesep \Tintlit(3)
\]

The corresponding derivation is:
\[
\begin{array}{rl}
\Nexpr & \derives\ \Nexpr \parsesep \Tplus \parsesep \underline{\Nexpr}\\
       & \derives\ \Nexpr \parsesep \Tplus \parsesep \Tintlit(3)\\
       & \derives\ \underline{\Nexpr} \parsesep \Tplus \parsesep \Tintlit(3)\\
       & \derives\ \Tintlit(1) \parsesep \Tplus \parsesep \Tintlit(3)\\
\end{array}
\]

This expression has the unique parse tree shown below.
\begin{center}
\begin{footnotesize}
\begin{adjustbox}{max width=\textwidth}
\Tree
[.$\Nexpr$
  [.$\Nexpr$ $\Tintlit(1)$ ]
  $\Tplus$
  [.$\Nexpr$ $\Tintlit(3)$ ]
]
\end{adjustbox}
\end{footnotesize}
\end{center}

\ExampleDef{Failing to Parse an Invalid Expression}
The expression \texttt{1 + * 3}
is transformed by the lexical analysis into:
\[
\Tintlit(1) \parsesep \Tplus \parsesep \Tmul \parsesep \Tintlit(3)
\]

An attempted derivation gets stuck as follows:
\[
\begin{array}{rl}
\Nexpr & \derives\ \Nexpr \parsesep \Tplus \parsesep \underline{\Nexpr}\\
       & \not\derives\ \Nexpr \parsesep \Tplus \parsesep \Tmul \parsesep \Tintlit(3)
\end{array}
\]
because there is no derivation for $\Nexpr$ whose first token is $\Tmul$.

\ExampleDef{Shift-Reduce Parsing with Binary Operators and Priority}
The expression \texttt{1 + 2 * 3}
is transformed by the lexical analysis into:
\[
\Tintlit(1) \parsesep \Tplus \parsesep \Tintlit(2) \parsesep \Tmul \parsesep \Tintlit(3)
\]

Since the derivation rules for expressions are ambiguous, this expression has
two possible parse trees:
\begin{center}
\begin{footnotesize}
\begin{tabular}{c|c}
\begin{adjustbox}{max width=0.45\textwidth}
\Tree
[.$\Nexpr$
  [.$\Nexpr$ $\Tintlit(1)$ ]
  $\Tplus$
  [.$\Nexpr$
    [.$\Nexpr$ $\Tintlit(2)$ ]
    $\Tmul$
    [.$\Nexpr$ $\Tintlit(3)$ ]
  ]
]
\end{adjustbox}
&
\begin{adjustbox}{max width=0.45\textwidth}
\Tree
[.$\Nexpr$
  [.$\Nexpr$
    [.$\Nexpr$ $\Tintlit(1)$ ]
    $\Tplus$
    [.$\Nexpr$ $\Tintlit(2)$ ]
  ]
  $\Tmul$
  [.$\Nexpr$ $\Tintlit(3)$ ]
]
\end{adjustbox}\\
\\
(a) intended parse tree.
&
(b) unintended parse tree.
\end{tabular}
\end{footnotesize}
\end{center}

Parse tree (a) is the intended parse tree, which reflects the fact
that $\Tmul$ has higher priority than $\Tplus$.

We show a run of a shift-reduce parser that produces the intended
parse tree (a) using the priority and associativity rules defined above.
%
Note that a shift-reduce parser stack is
represented by a box, where the bottom of the stack is on the left
and the top of the stack is on the right, and the input is represented by
a list of tokens.
%
An empty stack is represented by \fbox{$\emptysentence$}.
%
A node is created for each reduction step
for the non-terminal being reduced, using the stack content
as the children of the node.

\begin{center}
\begin{footnotesize}
\begin{tabular}{rll}
\textbf{\#} & \fbox{\textbf{Stack}} \textbf{Input} & \textbf{Action}\\
1 & \fbox{$\emptysentence$} $\Tintlit(1) \parsesep \Tplus \parsesep \Tintlit(2) \parsesep \Tmul \parsesep \Tintlit(3)$ & shift\\
2 & \fbox{$\Tintlit(1)$} $\Tplus \parsesep \Tintlit(2) \parsesep \Tmul \parsesep \Tintlit(3)$ & reduce $\Nexpr$\\
3 & \fbox{$\Nexpr$} $\Tplus \parsesep \Tintlit(2) \parsesep \Tmul \parsesep \Tintlit(3)$ & shift\\
4 & \fbox{$\Nexpr \parsesep \Tplus$} $\Tintlit(2) \parsesep \Tmul \parsesep \Tintlit(3)$ & shift\\
5 & \fbox{$\Nexpr \parsesep \Tplus \parsesep \Tintlit(2)$} $\Tmul \parsesep \Tintlit(3)$ & reduce $\Nexpr$\\
6 & \fbox{$\Nexpr \parsesep \Tplus \parsesep \Nexpr$} $\Tmul \parsesep \Tintlit(3)$ & shift\\
7 & \fbox{$\Nexpr \parsesep \Tplus \parsesep \Nexpr \parsesep \Tmul$} $\Tintlit(3)$ & shift\\
8 & \fbox{$\Nexpr \parsesep \Tplus \parsesep \Nexpr \parsesep \Tmul \parsesep \Tintlit(3)$} $\emptysentence$ & reduce $\Nexpr$\\
9 & \fbox{$\Nexpr \parsesep \Tplus \parsesep \Nexpr \parsesep \Tmul \parsesep \Nexpr$} $\emptysentence$ & reduce $\Nexpr$\\
10 & \fbox{$\Nexpr \parsesep \Tplus \parsesep \Nexpr$} $\emptysentence$ & reduce $\Nexpr$\\
11 & \fbox{$\Nexpr$} $\emptysentence$ & accept
\end{tabular}
\end{footnotesize}
\end{center}

\paragraph{Observation:}
Notice that there is no reduction to $\Nbinop$ in the table above.
This is because $\Nbinop$ is defined by an inlined derivation, so after expanding the grammar,
we have the grammar rules
$\Nexpr \derives \Nexpr \parsesep \Tplus \parsesep \Nexpr$ and
$\Nexpr \derives \Nexpr \parsesep \Tmul \parsesep \Nexpr$.

\paragraph{Observation:}
Step 6 shows how the parser resolves the ambiguity in\\
$\Tintlit(1) \parsesep \Tplus \parsesep \Tintlit(2) \parsesep \Tmul \parsesep \Tintlit(3)$.
At that point the stack contains
$\Nexpr \parsesep \Tplus \parsesep \Nexpr$
and the next token is $\Tmul$.
The parser shifts $\Tmul$ instead of reducing
$\Nexpr \parsesep \Tplus \parsesep \Nexpr$ immediately,
thereby allowing the higher-priority multiplication subexpression
$\Tintlit(2) \parsesep \Tmul \parsesep \Tintlit(3)$ to be formed first.
This is how the priority of $\Tmul$ over $\Tplus$ is enforced.

The table above corresponds to the following derivation in the expanded grammar
(the underline indicates the non-terminal that is reduced at each step):
\[
\begin{array}{rl}
\Nexpr & \derives\ \Nexpr \parsesep \Tplus \parsesep \underline{\Nexpr}\\
       & \derives\ \Nexpr \parsesep \Tplus \parsesep \Nexpr \parsesep \Tmul \parsesep \underline{\Nexpr}\\
       & \derives\ \Nexpr \parsesep \Tplus \parsesep \Nexpr \parsesep \Tmul \parsesep \Tintlit(3)\\
       & \derives\ \Nexpr \parsesep \Tplus \parsesep \underline{\Nexpr} \parsesep \Tmul \parsesep \Tintlit(3)\\
       & \derives\ \Nexpr \parsesep \Tplus \parsesep \Tintlit(2) \parsesep \Tmul \parsesep \Tintlit(3)\\
       & \derives\ \underline{\Nexpr} \parsesep \Tplus \parsesep \Tintlit(2) \parsesep \Tmul \parsesep \Tintlit(3)\\
       & \derives\ \Tintlit(1) \parsesep \Tplus \parsesep \Tintlit(2) \parsesep \Tmul \parsesep \Tintlit(3)\\
\end{array}
\]
Notice that the derivation corresponds perfectly to the concatenation of the stack
and input, read bottom-up in the table above, with the rightmost non-terminal
being the one that is reduced at each step.

\ExampleDef{Shift-Reduce Parsing with Left Associativity}
The expression \texttt{1 + 2 + 3}
is transformed by the lexical analysis into:
\[
\Tintlit(1) \parsesep \Tplus \parsesep \Tintlit(2) \parsesep \Tplus \parsesep \Tintlit(3)
\]

A corresponding shift-reduce parsing run is as follows:
\begin{center}
\begin{small}
\begin{tabular}{rll}
\textbf{Step} & \fbox{\textbf{Stack}} \textbf{Input} & \textbf{Action}\\
1 & \fbox{$\emptysentence$} $\Tintlit(1) \parsesep \Tplus \parsesep \Tintlit(2) \parsesep \Tplus \parsesep \Tintlit(3)$ & shift\\
2 & \fbox{$\Tintlit(1)$} $\Tplus \parsesep \Tintlit(2) \parsesep \Tplus \parsesep \Tintlit(3)$ & reduce $\Nexpr$\\
3 & \fbox{$\Nexpr$} $\Tplus \parsesep \Tintlit(2) \parsesep \Tplus \parsesep \Tintlit(3)$ & shift\\
4 & \fbox{$\Nexpr \parsesep \Tplus$} $\Tintlit(2) \parsesep \Tplus \parsesep \Tintlit(3)$ & shift\\
5 & \fbox{$\Nexpr \parsesep \Tplus \parsesep \Tintlit(2)$} $\Tplus \parsesep \Tintlit(3)$ & reduce $\Nexpr$\\
6 & \fbox{$\Nexpr \parsesep \Tplus \parsesep \Nexpr$} $\Tplus \parsesep \Tintlit(3)$ & reduce $\Nexpr$\\
7 & \fbox{$\Nexpr$} $\Tplus \parsesep \Tintlit(3)$ & shift\\
8 & \fbox{$\Nexpr \parsesep \Tplus$} $\Tintlit(3)$ & shift\\
9 & \fbox{$\Nexpr \parsesep \Tplus \parsesep \Tintlit(3)$} $\emptysentence$ & reduce $\Nexpr$\\
10 & \fbox{$\Nexpr \parsesep \Tplus \parsesep \Nexpr$} $\emptysentence$ & reduce $\Nexpr$\\
11 & \fbox{$\Nexpr$} $\emptysentence$ & accept
\end{tabular}
\end{small}
\end{center}

\paragraph{Observation:}
As explained in the previous example, there is no reduction to $\Nbinop$,
since $\Nbinop$ is an inlined non-terminal.
At step 6 the stack contains
$\Nexpr \parsesep \Tplus \parsesep \Nexpr$
and the next token is again $\Tplus$.
The parser reduces
$\Nexpr \parsesep \Tplus \parsesep \Nexpr$
to $\Nexpr$ instead of shifting the second $\Tplus$.
This is how the ambiguity is resolved using the fact that $\Tplus$
is declared $\leftassoc$:
when the competing actions have the same priority, left associativity chooses
reduce rather than shift.
As a result, the input is parsed as $(1 + 2) + 3$, not as $1 + (2 + 3)$.

The table above corresponds to the following derivation in the expanded grammar
(the underline indicates the non-terminal that is reduced at each step):
\[
\begin{array}{rl}
\Nexpr & \derives\ \Nexpr \parsesep \Tplus \parsesep \underline{\Nexpr}\\
       & \derives\ \Nexpr \parsesep \Tplus \parsesep \Tintlit(3)\\
       & \derives\ \underline{\Nexpr} \parsesep \Tplus \parsesep \Tintlit(3)\\
       & \derives\ \Nexpr \parsesep \Tplus \parsesep \underline{\Nexpr} \parsesep \Tplus \parsesep \Tintlit(3)\\
       & \derives\ \Nexpr \parsesep \Tplus \parsesep \Tintlit(2) \parsesep \Tplus \parsesep \Tintlit(3)\\
       & \derives\ \underline{\Nexpr} \parsesep \Tplus \parsesep \Tintlit(2) \parsesep \Tplus \parsesep \Tintlit(3)\\
       & \derives\ \Tintlit(1) \parsesep \Tplus \parsesep \Tintlit(2) \parsesep \Tplus \parsesep \Tintlit(3)\\
\end{array}
\]

\ExampleDef{Shift-Reduce Parsing with Unary Operators}
The expression \texttt{- - 1}
is transformed by the lexical analysis into:
\[
\Tminus \parsesep \Tminus \parsesep \Tintlit(1)
\]

The unique parse tree for this expression in the expanded grammar is:
\begin{center}
\begin{footnotesize}
\begin{adjustbox}{max width=\textwidth}
\Tree
[.$\Nexpr$
  $\Tminus$
  [.$\Nexpr$
    $\Tminus$
    [.$\Nexpr$ $\Tintlit(1)$ ]
  ]
]
\end{adjustbox}
\end{footnotesize}
\end{center}

A corresponding shift-reduce parsing run is as follows:
\begin{center}
\begin{small}
\begin{tabular}{rll}
\textbf{Step} & \fbox{\textbf{Stack}} \textbf{Input} & \textbf{Action}\\
1 & \fbox{$\emptysentence$} $\Tminus \parsesep \Tminus \parsesep \Tintlit(1)$ & shift\\
2 & \fbox{$\Tminus$} $\Tminus \parsesep \Tintlit(1)$ & shift\\
3 & \fbox{$\Tminus \parsesep \Tminus$} $\Tintlit(1)$ & shift\\
4 & \fbox{$\Tminus \parsesep \Tminus \parsesep \Tintlit(1)$} $\emptysentence$ & reduce $\Nexpr$\\
5 & \fbox{$\Tminus \parsesep \Tminus \parsesep \Nexpr$} $\emptysentence$ & reduce $\Nexpr$\\
6 & \fbox{$\Tminus \parsesep \Nexpr$} $\emptysentence$ & reduce $\Nexpr$\\
7 & \fbox{$\Nexpr$} $\emptysentence$ & accept
\end{tabular}
\end{small}
\end{center}

As explained in the previous examples, there is no reduction to $\Nunop$,
since it is an inlined non-terminal.

\paragraph{Observation:}
This example shows the parser building a nested unary expression.
After step 4, the stack suffix $\Tminus \parsesep \Tminus \parsesep \Nexpr$
matches the right-hand side obtained by expanding two applications of the unary-minus production.
The parser therefore performs two successive reductions of the form
$\Tminus \parsesep \Nexpr$ to $\Nexpr$,
which yields the parse $-(-1)$.

The table above corresponds to the following derivation in the expanded grammar
(the underline indicates the non-terminal that is reduced at each step):
\[
\begin{array}{rl}
\Nexpr & \derives\ \Tminus \parsesep \underline{\Nexpr}\\
       & \derives\ \Tminus \parsesep \Tminus \parsesep \underline{\Nexpr}\\
       & \derives\ \Tminus \parsesep \Tminus \parsesep \Tintlit(1)\\
\end{array}
\]


\ExampleDef{Shift-Reduce Parsing with Unary Operators and Priority}
The expression \texttt{-1 - 2}
is transformed by the lexical analysis into:
\[
\Tminus \parsesep \Tintlit(1) \parsesep \Tminus \parsesep \Tintlit(2)
\]

A corresponding pair of parse trees for the expanded grammar is:
\begin{center}
\begin{footnotesize}
\begin{tabular}{c|c}
\begin{adjustbox}{max width=0.45\textwidth}
\Tree
[.$\Nexpr$
  [.$\Nexpr$
    $\Tminus$
    [.$\Nexpr$ $\Tintlit(1)$ ]
  ]
  $\Tminus$
  [.$\Nexpr$ $\Tintlit(2)$ ]
]
\end{adjustbox}
&
\begin{adjustbox}{max width=0.45\textwidth}
\Tree
[.$\Nexpr$
  $\Tminus$
  [.$\Nexpr$
    [.$\Nexpr$ $\Tintlit(1)$ ]
    $\Tminus$
    [.$\Nexpr$ $\Tintlit(2)$ ]
  ]
]
\end{adjustbox}
\\
(a) intended parse tree. &
(b) unintended parse tree.
\end{tabular}
\end{footnotesize}
\end{center}

A corresponding shift-reduce parsing run that produces the intended parse tree
is as follows:
\begin{center}
\begin{small}
\begin{tabular}{rll}
\textbf{Step} & \fbox{\textbf{Stack}} \textbf{Input} & \textbf{Action}\\
1 & \fbox{$\emptysentence$} $\Tminus \parsesep \Tintlit(1) \parsesep \Tminus \parsesep \Tintlit(2)$ & shift\\
2 & \fbox{$\Tminus$} $\Tintlit(1) \parsesep \Tminus \parsesep \Tintlit(2)$ & shift\\
3 & \fbox{$\Tminus \parsesep \Tintlit(1)$} $\Tminus \parsesep \Tintlit(2)$ & reduce $\Nexpr$\\
4 & \fbox{$\Tminus \parsesep \Nexpr$} $\Tminus \parsesep \Tintlit(2)$ & reduce $\Nexpr$\\
5 & \fbox{$\Nexpr$} $\Tminus \parsesep \Tintlit(2)$ & shift\\
6 & \fbox{$\Nexpr \parsesep \Tminus$} $\Tintlit(2)$ & shift\\
7 & \fbox{$\Nexpr \parsesep \Tminus \parsesep \Tintlit(2)$} $\emptysentence$ & reduce $\Nexpr$\\
8 & \fbox{$\Nexpr \parsesep \Tminus \parsesep \Nexpr$} $\emptysentence$ & reduce $\Nexpr$\\
9 & \fbox{$\Nexpr$} $\emptysentence$ & accept
\end{tabular}
\end{small}
\end{center}

As explained in the previous example, there are no reductions to $\Nunop$ or $\Nbinop$,
since they are inlined non-terminals.

\paragraph{Observation:}
At step 4, stack contains $\Tminus \parsesep \Nexpr$ and the next token is again $\Tminus$.
The parser reduces $\Tminus \parsesep \Nexpr$ to $\Nexpr$ instead of shifting the second $\Tminus$.
This is exactly the priority rule for unary minus:
the production $\Nexpr \derives \Tminus \parsesep \Nexpr$ has priority $\Tunops$,
which is higher than the binary-operator priority of $\Tminus$.
As a result, the input is parsed as $(-1) - 2$, not as $-(1 - 2)$.

The table above corresponds to the following derivation in the expanded grammar
(the underline indicates the non-terminal that is reduced at each step):
\[
\begin{array}{rl}
\Nexpr & \derives\ \underline{\Nexpr} \parsesep \Tminus \parsesep \Nexpr\\
       & \derives\ \underline{\Nexpr} \parsesep \Tminus \parsesep \Tintlit(2)\\
       & \derives\ \Tminus \parsesep \underline{\Nexpr} \parsesep \Tminus \parsesep \Tintlit(2)\\
       & \derives\ \Tminus \parsesep \Tintlit(1) \parsesep \Tminus \parsesep \Tintlit(2)\\
\end{array}
\]

\ExampleDef{Parsing with Parametric Productions}
The expression \texttt{b[c,d,e]}
is transformed by the lexical analysis into:
\[
\Tidentifier(b) \parsesep \Tlbracket \parsesep \Tidentifier(c) \parsesep \Tcomma \parsesep \Tidentifier(d) \parsesep \Tcomma \parsesep \Tidentifier(e) \parsesep \Trbracket
\]

\texthypertarget{def-clistoneslice}
To emphasize the role of parametric productions, let $\clistoneslice$
be the unique symbol introduced when expanding $\ClistOne{\Nslice}$.
The relevant expansion is
\[
\clistoneslice \derives \Nslice \;|\; \Nslice \parsesep \Tcomma \parsesep \clistoneslice
\]
and the derivation for $\Nslices$ becomes
\[
\Nslices \derives \Tlbracket \parsesep \clistoneslice \parsesep \Trbracket \enspace.
\]
